\chapter{A unified framework for finite-dimensional functional theories}
\label{chap:unified}

% D4 statements written 2026-11-26 after C.-C. Wang, J. Liebert, M. Penz, C. Schilling, arXiv:2606.06676,
% and approved by Miguel 2026-11-26. Not yet formalized (no \lean references until the declarations exist).

All functional theories on a finite-dimensional Hilbert space share one structure: a space of potentials coupled
linearly to a fixed Hamiltonian. The results of this chapter are proved once for that structure; lattice DFT
(Chapter~\ref{chap:ks}), RDMFT and spin DFT are instances. Throughout, $\mathcal H=\mathbb C^d$, states are unit vectors
$\psi$ or density matrices $\Gamma\succeq0$ with $\operatorname{tr}\Gamma=1$.

\section{Scopes, densities and functionals}

\begin{definition}[Scope]\label{def:scope}
  A scope is $S=(V,\iota,H_0)$: a finite-dimensional real vector space $V$, a linear map $\iota:V\to\mathbb C^{d\times d}$ with
  Hermitian values, and a Hermitian $H_0$. The Hamiltonians are $H(v)=H_0+\iota(v)$. The scope is \emph{abelian} if the
  $\iota(v)$ commute, and \emph{non-redundant} if $\iota$ is injective and $1\notin\iota(V)$.
\end{definition}

\begin{definition}[Densities, ranges, functionals]\label{def:scope-functionals}
  \uses{def:scope}
  The density of $\Gamma$ is $\rho_\Gamma\in V^*$, $\rho_\Gamma(v)=\operatorname{Re}\operatorname{tr}(\Gamma\iota(v))$; for $\Gamma=|\psi\rangle\langle\psi|$ write
  $\rho_\psi$. The ranges are $B_p=\{\rho_\psi\}$ and $B_e=\{\rho_\Gamma\}$. The ground-state energy is
  $E(v)=\inf_\psi\langle\psi,H(v)\psi\rangle$. The constrained-search functionals are
  $F_p(\rho)=\inf\{\langle\psi,H_0\psi\rangle:\rho_\psi=\rho\}$ and $F_e(\rho)=\inf\{\operatorname{tr}(H_0\Gamma):\rho_\Gamma=\rho\}$, with value $+\infty$
  off $B_p$, $B_e$.
\end{definition}

\begin{definition}[Critical and regular values]\label{def:scope-critical}
  \uses{def:scope-functionals}
  A unit vector $\psi$ is critical if it is an eigenvector of some $\iota(v)\notin\mathbb R\cdot1$. $B_c$ is the set of densities of
  critical states and $B_r=B_p\setminus B_c$ the set of regular values.
\end{definition}

\section{Geometry and duality (stage 1)}

\begin{theorem}[Observable ranges]\label{thm:scope-ranges}
  \uses{def:scope-functionals}
  $B_p$ and $B_e$ are compact and $B_e=\operatorname{conv}B_p$. If the scope is abelian, $B_p=B_e$ is the convex hull of the finitely
  many joint-eigenvalue densities.
\end{theorem}

\begin{theorem}[Constrained-search functionals]\label{thm:scope-functionals}
  \uses{def:scope-functionals, thm:scope-ranges}
  The infima defining $F_p$, $F_e$ are attained on the ranges. $E(v)=\inf_\rho(F_p(\rho)+\rho(v))=\inf_\rho(F_e(\rho)+\rho(v))$.
  $F_e$ is convex, $F_p$ and $F_e$ are lower semicontinuous, $F_e\le F_p$, $F_e(\rho)=\sup_v(E(v)-\rho(v))$ on $B_e$, and $F_e$ is
  the greatest convex lower semicontinuous function below $F_p$.
\end{theorem}

\begin{theorem}[Critical states]\label{thm:scope-critical}
  \uses{def:scope-critical}
  For a non-redundant scope with basis $B_i=\iota(e_i)$, $\psi$ is critical iff $\psi,B_1\psi,\dots,B_m\psi$ are linearly dependent
  over $\mathbb R$, iff their real Gram determinant vanishes.
\end{theorem}

\begin{theorem}[Hohenberg--Kohn]\label{thm:scope-hk}
  \uses{thm:scope-critical}
  (Weak.) If ground states of $H(v_1)$ and $H(v_2)$ have the same density, each is a ground state of the other Hamiltonian;
  likewise for ensembles. (Strong.) If $\rho\in B_r$ is the density of a pure ground state of $H(v)$ and of $H(v')$, then
  $\iota(v-v')\in\mathbb R\cdot1$; for a non-redundant scope, $v=v'$.
\end{theorem}

\begin{theorem}[Ensemble $v$-representability]\label{thm:scope-ens-vrep}
  \uses{thm:scope-functionals}
  Every $\rho$ in the intrinsic interior of $B_e$ is the density of an ensemble ground state of some $H(v)$, with $-v$ a
  subgradient of $F_e$ at $\rho$. If $\rho$ is the density of a pure ground state, $F_p(\rho)=F_e(\rho)$.
\end{theorem}

\begin{theorem}[Instances]\label{thm:scope-instances}
  \uses{thm:scope-ranges, thm:scope-functionals, thm:lieb_duality, thm:coleman}
  Lattice DFT ($\iota(v)=\sum_iv_in_i$ on $N$ fermions): $B_p=B_e$ is the $(M,N)$-hypersimplex. RDMFT ($\iota$ onto all one-body
  operators): $B_e=\{0\le\gamma\le1,\operatorname{tr}\gamma=N\}$ (Coleman). Spin chains ($\iota(e_i)=Z_i$): $B_p=B_e=[-1,1]^m$. The lattice
  Lieb duality is the instance of Theorem~\ref{thm:scope-functionals}.
\end{theorem}

\section{Representability by eigenstates (stage 2)}

\begin{theorem}[Eigenstate $v$-representability]\label{thm:scope-eig-vrep}
  \uses{thm:scope-functionals, def:scope-critical}
  Every $\rho\in B_r$ is $\rho_\psi$ for a unit vector $\psi$ with $H(v)\psi=\lambda\psi$ for some $v\in V$, $\lambda\in\mathbb R$, and $\psi$ attains
  $F_p(\rho)$. (To check during the proof: whether non-redundancy is needed.) $F_e$ is continuous on the intrinsic interior
  of $B_e$.
\end{theorem}

\section{Change of scope and examples (stage 3)}

\begin{theorem}[Reduction and purification]\label{thm:scope-change}
  \uses{def:scope-functionals}
  For linear $f:V_r\to V$ and $v_0\in V$, the reduced scope $(V_r,\iota\circ f,H_0+\iota(v_0))$ has densities $f^*\circ\rho$.
  $F_e$ of a scope equals $F_p$ of its purified scope on $\mathcal H\otimes\mathcal H$ (exact form to be checked against
  arXiv:2606.06676, \S VII.B).
\end{theorem}

\begin{theorem}[Examples]\label{thm:scope-examples}
  \uses{def:scope-functionals}
  For $\mathcal H=\mathbb C^2$ and $\iota(v)=v_1X+v_2Y+v_3Z$, $B_p$ is the unit sphere and $B_e$ the unit ball, so $B_p$ need not be
  convex. For two qubits with observables $X_1,Y_1,Z_1,Z_2$ (Appendix A of the paper) there is an interior point with
  $F_e<F_p$ (if the values are exact).
\end{theorem}
